\centerline{\bf \Huge Муниципальная геома}

\problem{1}
$\bigl[\textbf{Кировская область, 24/25, 9.4.}\bigr]$
В треугольник $АВС$ вписана окружность, а в нее --- треугольник $DEF$ со сторонами, параллельными сторонам треугольника $АВС$. Верно ли, что отношение площади треугольника $D E F$ к площади треугольника $A B C$ обязательно больше 1/2024?

\problem{2}
$\bigl[\textbf{Пермь, 24/25, 9.3.}\bigr]$
В остроугольном треугольнике $A B C$ провели биссектрису $A F$, высоту $B H$ и медиану $C N$. Оказалось, что $F A$ - биссектриса угла $N F H$, а прямые $H B$ и $N F$ перпендикулярны. Докажите, что треугольник $A B C$ является равносторонним.


\problem{3}
$\bigl[\textbf{Амурская область, 23/24, 10.3.}\bigr]$
Из произвольной точки $M$ внутри данного угла $A$ опущены перпендикуляры $M P$ и $M Q$ на стороны угла. Из точки $A$ перпендикуляр $A K$ на отрезок $P Q$. Докажите, что $\angle P A K=\angle M A Q$.

\problem{4}
$\bigl[\textbf{Московская обалсть, 24/25, 10.5.}\bigr]$
В трапеции $A B C D$ длина боковой стороны $A B$ равна сумме длин оснований. Окружность $\Omega$ с центром в точке $A$, проходящая через $D$, пересекает отрезки $B D$ и $B A$ в точках $E$ и $F$ соответственно. Докажите, что точки $B, C, E, F$ лежат на одной окружности.

\problem{5}
$\bigl[\textbf{Республика Башкортостан, 23/24, 11.5.}\bigr]$
Окружность, проходящая через вершины $A$ и $B$ треугольника $A B C$ пересекает его стороны $B C$ и $A C$ во внутренних точках $Q$ и $P$ соответственно, причем $A P=3 P C$. Докажите, что $\angle A B M=\angle M Q P$, где $B M$ --- медиана треугольника $A B C$.

